\section{The Verifier as a Process Reward}
\label{sec:reward}

Runtime governance (\S\ref{sec:eval}) keeps the verifier in the loop at inference. We now ask whether the same product-order geometry can be \emph{internalized}: used as a training signal, does it teach the policy to recover on its own, so that recovery survives when the gate is later removed? We reuse the identical verifier---its graded verdict and the persisted branch state $\Phi$ of Eq.~\ref{eq:phi}---as a process reward for GRPO~\cite{deepseekmath2024grpo,lightman2024prm}, with no domain-specific reward engineering.

\paragraph{Setup.}
We fine-tune Qwen3-8B with LoRA on the $9$ sub-saturated multi-action ANM scenarios (greedy decoding, thinking off). Three reward arms share an identical rollout gate (\texttt{progress\_mag}) and differ only in the scalar reward, isolating the reward signal: \textbf{D} (\emph{geometric}) rewards descent in the product order over $\Phi=(S,\sigma)$; \textbf{E} (\emph{scalar}) rewards the count-delta of cleared violations; and an ablation \textbf{D-flat} strips the geometric component from D, leaving only the verdict ladder. The operating point is iter\_2 (fixed in advance; iter\,1 undertrained, iter\,3+ over-optimized). Our primary outcome is \emph{ungated} recovery---the gate is removed at evaluation, so recovery measures what the policy internalized rather than what the gate enforces; we report it pooled over $5$ training seeds with Wilson 95\% CI, alongside gated recovery and steps-to-recover.

\begin{table}[t]
\centering
\caption{Verifier-as-reward on the $9$ sub-saturated ANM scenarios ($5$ training seeds pooled, fixed iter\_2; recovery [Wilson 95\% CI]). Ungated recovery (gate removed at eval) measures internalization: only the full geometric reward (D) holds at or above the untrained base; stripping the geometry (D-flat) or reducing it to a scalar count (E) regresses below it. The untrained base runs under the same \texttt{progress\_mag} gate (no LoRA), so its gated and ungated rates coincide. Steps-to-recover is the mean over recovered episodes.}
\label{tab:reward}
\begin{tabular}{llccc}
\toprule
Arm & reward signal & gated & \textbf{ungated} & steps$\downarrow$ \\
\midrule
base (untrained) & --- & $0.778$ & $0.778$ & $4.71$ \\
\textbf{D} & geometric ($\Phi$ descent) & $\mathbf{0.956}$ [.85,.99] & $\mathbf{0.844}$ [.71,.92] & $\mathbf{3.64}$ \\
D-flat & verdict ladder only & $0.889$ & $0.733$ [.59,.84] & $4.08$ \\
E & scalar count-delta & $0.889$ & $0.667$ [.52,.79] & $3.84$ \\
\bottomrule
\end{tabular}
\end{table}

\paragraph{Geometry is the active ingredient.}
Table~\ref{tab:reward} reads as a monotone ablation chain, $\text{D} > \text{D-flat} > \text{E}$. The full geometric reward is the only arm that does \emph{not} regress below the untrained base once the gate is removed ($0.844$, base $0.778$); stripping the geometry (D-flat, $0.733$) or replacing it with a scalar count (E, $0.667$) leaves the policy \emph{worse} than untrained---a reward that does not encode the geometry teaches reliance on the gate rather than internalized recovery. The geometric component accounts for $+0.11$ of this separation and the verdict grading a further $+0.066$; D also recovers fastest ($3.64$ vs.\ base $4.71$ steps). The in-distribution separation is directional---the arm CIs overlap at $N{=}9$---so we do not read it as a significant ungated lift over base; the held-out test below carries the statistical weight. We likewise do not treat the gated D-vs-E gap as evidence of reward quality: the rollout gate enforces the same product-order geometry that D's reward encodes, making that comparison partially self-referential.

\paragraph{Internalization generalizes.}
On $15$ held-out scenarios disjoint from training, D recovers $14/15$ ungated against the untrained base's $9/15$ (one-sided Fisher exact $p{=}0.04$)---the policy carries the geometry to unseen instances, ruling out in-distribution memorization. The scalar arm E reaches the same ungated rate ($14/15$), but its gated recovery falls to $0.80$, below base's $0.87$: the two arms separate by \emph{mechanism} rather than headline rate---D degrades gracefully when the gate is removed, whereas E overfits to the gate's support filter. (Held-out figures are a single adapter, seed $0$.)

\paragraph{Scope.}
This is a single-violation-family, single-model ($8$B) study; the in-distribution effect is modest relative to evaluation noise---identical greedy decodes vary by $\pm1/9$ across runs (GPU matmul, temperature $0$), which is why we pool five training seeds and anchor on the held-out test. Concurrently activating multiple, physically incomparable violation families---where a scalar surrogate must collapse the product order and the geometric advantage should be largest---is the natural next test (\S\ref{sec:discussion}).
